\documentclass{situationist-preprint}

\title{Petitionary Prayer as Non-Causal Difference-Making}
\msid{REL-12-a}
\msfamily{REL-12}
\msdiscipline{Philosophy of Religion}
\mstype{Principal manuscript}
\mskeywords{petitionary prayer; answered prayer; divine reasons; overdetermination; counterfactual dependence; timelessness}

\begin{document}
\maketitle

\begin{abstract}
If God is perfectly good, omniscient and omnipotent, God will bring about whatever is best whether or not anyone asks; petitionary prayer, it seems, can make no difference. The usual replies argue that some goods are available only if requested, so that the request makes a causal or counterfactual difference to the outcome. We argue for a different answer. An outcome is granted in response to a request when the request is among the reasons for which God brings it about, and this can be so even when God would have brought about the same outcome had there been no request. The request then makes no difference to whether the outcome occurs, and is not among its causes, on a counterfactual account of causation. It makes a difference to what the outcome is: a granting, the answer to a particular person's asking. We show that this relation is coherent for a timeless God, that it does not reduce to the dependence of the outcome on the request, and that it explains why answered prayer and unanswered prayer can coincide in every observable respect.
\end{abstract}

\section{The problem}

Petitionary prayer asks God to bring something about: the recovery of a friend, rain for a harvest, peace. The practice presupposes that asking can matter. A familiar argument suggests that it cannot. A perfectly good God brings about an outcome if and only if it is best that it obtain; an omniscient God knows whether it is best; an omnipotent God can bring it about. Whether the outcome obtains is therefore settled by whether it is best, and a request does not change what is best \citep{Basinger1983}.

Two families of reply have dominated. The first holds that a request can change what is best: some goods, such as the dependence of creatures on God and a relationship of asking and giving, are realized only if God sometimes waits to be asked, so that God has reason to make some outcomes depend on requests \citep{Stump1979,MurrayMeyers1994}. On this view prayer makes a counterfactual difference: in some cases, had the request not been made, the outcome would not have occurred. The second family doubts that the conditions for answered prayer can be satisfied or known. \citet{Davison2017} argues that a prayer is answered only if God brings about the outcome at least partly because of the prayer, and that it is very hard to show that this condition is ever met, since a perfectly good God might have had sufficient reason to bring about the outcome in any case.

We accept Davison's condition and reject the inference. That God would have brought about an outcome anyway does not show that God did not bring it about because of the prayer. Reasons can overdetermine an action: an agent may act for two reasons, each sufficient. When God's reasons for an outcome include the request, the outcome is granted, whether or not it would have occurred without the request. Prayer then makes a difference without causal efficacy: an outcome may be granted because it was requested, even when it would otherwise have occurred.

\section{Framework}

\begin{definition}[Acting for a reason]\label{d:for}
An agent brings about $O$ \emph{for} a reason $r$ if $r$ is among the considerations in the light of which the agent wills $O$, so that a correct answer to the question why the agent willed $O$ cites $r$.
\end{definition}

\begin{definition}[Granting]\label{d:grant}
An outcome $O$ is \emph{granted} in response to a request $q$ if God brings about $O$ and $q$ is among the reasons for which God brings it about.
\end{definition}

\begin{definition}[Causal difference]\label{d:causal}
A request $q$ makes a \emph{causal difference} to $O$ if $O$ counterfactually depends on $q$: had $q$ not been made, $O$ would not have occurred.
\end{definition}

Definition~\ref{d:causal} follows the account of causation of MET-04, on which a cause is a counterfactual difference-maker selected by norms of normality. A request that is not a counterfactual difference-maker for $O$ is not, on that account, among the causes of $O$.

\begin{principle}[Overdetermination of reasons]\label{p:over}
An agent may act for two reasons each of which would by itself have been sufficient for the action, and act for both.
\end{principle}

Principle~\ref{p:over} is familiar from human agency. A doctor treats a patient because the patient asked and because the treatment is indicated; she would have proposed it had the patient not asked, and she treats the patient for both reasons. That the patient's asking was not necessary does not make it false that she acted on it. An explanation that cites an agent's reason says what she acted on; it does not say that, without that reason, she would not have acted. What depends counterfactually on the patient's asking is not the treatment but the treatment's being given at her request.

\section{The argument}

\begin{lemma}\label{l:indep}
Granting does not entail causal difference: $O$ may be granted in response to $q$ although $O$ would have occurred had $q$ not been made.
\end{lemma}

\begin{proof}
Suppose God has two reasons for bringing about a friend's recovery: that the recovery is good for the friend, and that the petitioner asked for it. Suppose the first is by itself sufficient, so that, had the petitioner not asked, God would have brought about the recovery for the first reason. By Principle~\ref{p:over}, God may nonetheless bring about the recovery for both. By Definition~\ref{d:grant}, the recovery is granted in response to the request. By Definition~\ref{d:causal}, the request makes no causal difference, since the recovery would have occurred without it.
\end{proof}

\begin{lemma}\label{l:diff}
Where $O$ is granted in response to $q$, the fact that $O$ is a granting counterfactually depends on $q$.
\end{lemma}

\begin{proof}
Had $q$ not been made, $q$ could not have been among God's reasons for $O$, since a request that is not made is not a consideration in the light of which anyone acts. So $O$ would not have been granted in response to $q$, even if $O$ had occurred.
\end{proof}

\begin{theorem}\label{t:main}
Petitionary prayer can make a difference without causal efficacy: an outcome may be granted because it was requested, even when it would otherwise have occurred.
\end{theorem}

\begin{proof}
By Lemma~\ref{l:indep}, an outcome may be granted in response to a request on which it does not counterfactually depend, so that the request is not among its causes. By Lemma~\ref{l:diff}, the request makes a difference to whether the outcome is a granting, which is a fact about the world: the outcome stands in a relation to the petitioner, as the answer to her asking, in which it would not otherwise have stood.
\end{proof}

The difference is not merely verbal. Whether a gift was given in answer to a request, or would have been given anyway and was given without regard to it, matters to the relation between giver and recipient, and to what gratitude is fitting. A child who asks for something her parent intended to give, and receives it as an answer, has been heard; a child whose request was ignored, though she received the same thing, has not. Theorem~\ref{t:main} locates the point of petitionary prayer in this relation, not in the extraction of outcomes.

\begin{proposition}[Timelessness]\label{p:timeless}
Granting is available to a timeless God.
\end{proposition}

\begin{proof}
REL-19 concludes that God is timeless. A timeless God cannot be causally affected by a temporal event, since causal influence runs from earlier to later. But Definition~\ref{d:grant} requires only that the request be among God's reasons, and God's timeless will can be responsive to temporal events in the way that, on the account of REL-05, God's timeless knowledge depends counterfactually on free choices: had the request not been made, God's will would not have included it among its reasons. Responsiveness of this kind is dependence of God's reasons on the request, not causation of God's willing by the request.
\end{proof}

Proposition~\ref{p:timeless} also disposes of a puzzle about prayers made after the outcome they ask for, such as a prayer that a letter already sent should contain good news. If the outcome is granted, the request is among God's timeless reasons for it, and nothing requires the request to precede the outcome; the prayer is not a cause, and so is not a backward cause. We do not claim that such prayers are appropriate when the petitioner knows the outcome; REL-22 discusses what can and cannot be done about the past.

\section{Objections and replies}

\begin{objection}[Idle reasons]
A reason that would make no difference to what an agent does is not one she acts on; it is at most one she approves of. If God would have brought about the recovery anyway, the request is idle.
\end{objection}

\begin{reply}
The objection identifies acting on a reason with the reason's being necessary for the action, and the identification is false for human agents, as the doctor's case shows. What distinguishes a reason acted on from one merely approved is its role in the agent's deliberation and willing, not its counterfactual indispensability. An omniscient agent who wills an outcome in view of everything that bears on it, the request included, acts on the request in the sense of Definition~\ref{d:for}.
\end{reply}

\begin{objection}[Measurement]
If answered prayer need not change outcomes, then nothing distinguishes a world in which prayers are answered from one in which they are not, and the claim that they are answered is empty.
\end{objection}

\begin{reply}
The two worlds differ in God's reasons, and so in the relations between outcomes and petitioners. They need not differ in the frequencies of outcomes, which is why controlled studies of intercessory prayer, such as \citet{Benson2006}, which found no benefit of intercessory prayer for recovery after cardiac surgery, test a hypothesis that Theorem~\ref{t:main} does not assert. The companion manuscript REL-12-b develops the point that the efficacy of prayer, in the sense here defended, cannot be measured. A claim about the reasons for which an agent acts is not empty because it does not show in outcome statistics.
\end{reply}

\begin{objection}[Causal difference after all]
On the Stump–Murray–Meyers view, some outcomes are good only if requested. Then requests sometimes make a causal difference, and Theorem~\ref{t:main} is unnecessary.
\end{objection}

\begin{reply}
We do not deny that requests sometimes make a causal difference; the conclusion is that they need not, and that prayer is answered in the cases where they do not. The Stump–Murray–Meyers view, if true, covers only outcomes that would not occur without a request. It leaves unexplained the most common case of petition, in which one asks for what a good God would plausibly do anyway, such as the relief of suffering. Theorem~\ref{t:main} covers that case.
\end{reply}

\section{Conclusion}

Petitionary prayer does make a difference, but not the difference the standard argument looks for. An outcome is granted when the request is among the reasons for which God brings it about, and this can be so when the outcome would have occurred without the request. The request is then not among the outcome's causes, and it is not a backward cause of anything; the difference it makes is that the outcome is an answer.

\begin{thebibliography}{99}\small
\bibitem[Basinger(1983)]{Basinger1983} Basinger, D. (1983). Why petition an omnipotent, omniscient, wholly good God? \emph{Religious Studies}, 19(1), 25--41.
\bibitem[Benson et~al.(2006)]{Benson2006} Benson, H. et~al. (2006). Study of the therapeutic effects of intercessory prayer (STEP) in cardiac bypass patients: A multicenter randomized trial of uncertainty and certainty of receiving intercessory prayer. \emph{American Heart Journal}, 151(4), 934--942.
\bibitem[Davison(2017)]{Davison2017} Davison, S.~A. (2017). \emph{Petitionary Prayer: A Philosophical Investigation}. Oxford University Press.
\bibitem[Murray and Meyers(1994)]{MurrayMeyers1994} Murray, M.~J. and Meyers, K. (1994). Ask and it will be given to you. \emph{Religious Studies}, 30(3), 311--330.
\bibitem[Stump(1979)]{Stump1979} Stump, E. (1979). Petitionary prayer. \emph{American Philosophical Quarterly}, 16(2), 81--91.
\end{thebibliography}

\end{document}
